\documentclass[preprint,3p]{elsarticle}

\usepackage{lineno,hyperref}
\usepackage{tabularx}
\usepackage{multirow}
\usepackage{listings}
\usepackage{color,soul}

%\usepackage{subfigure}

\usepackage{graphicx}
\usepackage{subcaption}

\modulolinenumbers[5]

\journal{Future Generation Computer Systems}

%\newcommand\boldred[1]{\textcolor{red}{\textbf{#1}}}
%\newcommand\boldblue[1]{\textcolor{blue}{\textbf{#1}}}
%\newcommand\boldcyan[1]{\textcolor{cyan}{\textbf{#1}}}

\newcommand\boldred[1]{#1}
\newcommand\boldblue[1]{#1}
\newcommand\boldcyan[1]{#1}


%%%%%%%%%%%%%%%%%%%%%%%
%% Elsevier bibliography styles
%%%%%%%%%%%%%%%%%%%%%%%
%% To change the style, put a % in front of the second line of the current style and
%% remove the % from the second line of the style you would like to use.
%%%%%%%%%%%%%%%%%%%%%%%

%% Numbered
%\bibliographystyle{model1-num-names}

%% Numbered without titles
%\bibliographystyle{model1a-num-names}

%% Harvard
%\bibliographystyle{model2-names.bst}\biboptions{authoryear}

%% Vancouver numbered
%\usepackage{numcompress}\bibliographystyle{model3-num-names}

%% Vancouver name/year
%\usepackage{numcompress}\bibliographystyle{model4-names}\biboptions{authoryear}

%% APA style
%\bibliographystyle{model5-names}\biboptions{authoryear}

%% AMA style
%\usepackage{numcompress}\bibliographystyle{model6-num-names}

%% `Elsevier LaTeX' style
\bibliographystyle{elsarticle-num}
%%%%%%%%%%%%%%%%%%%%%%%


\begin{document}

\begin{frontmatter}

\title{
	Blockchain-based resource accounting for a serverless execution platform
}

%% Group authors per affiliation:
\author{Jonatan Enes\corref{cor}}
\ead{jonatan.enes@udc.es}
\author{Roberto R. Exp\'osito\corref{cor2}}
\ead{rreye@udc.es}
\author{Juan Touri\~no\corref{cor2}}
\ead{juan@udc.es}

%[gac]
\address{Universidade da Coru\~na, CITIC, Computer Architecture Group, Campus de A Coru\~na, Spain}

\cortext[cor]{Corresponding author}

\begin{abstract}

\end{abstract}

\begin{keyword}
anonymous computing, decentralized accounting
\end{keyword}

\end{frontmatter}

\linenumbers

\section{Introduction}
\label{sec:intro}

\section{Experiments}
\label{sec:experiments}

\section{Común a todos los experimentos}

\begin{itemize}
	\item Se usaría 1 solo contenedor, dentro de 1 app y asociada a 1 usuario. El objetivo del artículo no sería demostrar tanto la plataforma serverless y sus capacidades de gestión de múltiples contenedores (eso ya lo han hecho otros artículos), sino sus funciones de contabilidad e integración con la blockchain.
\end{itemize}

\newcommand{\ContainerUp}{Se levanta un contenedor sin nada ejecutándose, con unos recursos asignados iniciales}

\newcommand{\UserSendsCredit}{El usuario manda crédito}
\newcommand{\UserNoCredit}{El usuario no tiene crédito}
\newcommand{\UserAccountingActivated}{Se activa el accounting para el usuario}


\newcommand{\ServerlessDown}{La plataforma (serverless) ajusta los recursos a la baja}
\newcommand{\ServerlessUp}{La plataforma (serverless) ajusta los recursos al alza}
\newcommand{\ServerlessAdapts}{La plataforma (serverless) ajusta los recursos según la demanda}

\newcommand{\AccountingIdle}{La plataforma (accounting) no detecta uso así que no consume nada}
\newcommand{\AccountingConsumes}{La plataforma (accounting) detecta uso y va consumiendo el crédito}
\newcommand{\AccountingRestricts}{La plataforma (accounting) restringe los recursos}
\newcommand{\AccountingAllows}{La plataforma (accounting) levanta la restricción}
\newcommand{\AccountingCredit}{El crédito restante debería ser}

\newcommand{\AppStart}{Arranca la ejecución de}
\newcommand{\AppRuns}{Continúa la ejecución de la aplicación}
\newcommand{\AppWaits}{La aplicación se para a la espera de que haya crédito}
\newcommand{\AppResume}{La aplicación se reanuda}


\newcommand{\Wait}{Se espera un tiempo}

\newpage
\section{Basic Experiment}
La secuencia sería:
\begin{itemize}
	\setlength\itemsep{0em}
	\item \ContainerUp
	\item \UserNoCredit
	\item \UserSendsCredit, poco ($\approx 6~GRC$)
	\item \UserAccountingActivated
	\item \ServerlessDown
	\item \AccountingIdle
	\item \AppStart~\textbf{SimpleApp}
	\item \ServerlessUp
	\item \AccountingConsumes
	\item \AppRuns~hasta un punto en el que el crédito restante sea bajo ($ < 2~GRC$)
	\item \UserSendsCredit, una cantidad justa ($ > 4~GRC$)
	\item \AppRuns~hasta que se consume todo el crédito 
	\item \AccountingRestricts
	\item \AppRuns~con los recursos restringidos (mínimos)
	\item \AccountingConsumes, se genera una ligera deuda
	\item \UserSendsCredit, una cantidad generosa ($ > 10~GRC$)
	\item \AccountingAllows
	\item \ServerlessAdapts
	\item \AppRuns~hasta su finalización
	\item \ServerlessDown
	\item \AccountingIdle
	\item \AccountingCredit~$ > 0~GRC$
\end{itemize}

Sobre la aplicación \textbf{SimpleApp}:
\begin{itemize}
	\setlength\itemsep{0em}
	\item Carga media (6-8 cores)
	\item Poco o nada sensible a cambios de recursos, esta aplicación va a tener que seguir ejecutándose aún bajo la restricción de recursos
	\item Con picos y valles de consumo, aunque no debería imponer un alto estrés a la plataforma en su gestión de recursos serverless
	\item Candidatos: 
	\begin{itemize}
		\item Benchmark de procesamiento numérico puro
	\end{itemize}
\end{itemize}

\newpage
\section{Trickle processing}
La secuencia sería:
\begin{itemize}
	\setlength\itemsep{0em}
	\item \ContainerUp
	\item \UserNoCredit
	\item \UserSendsCredit, muy poco ($\approx 1~GRC$)
	\item \UserAccountingActivated
	\item \ServerlessDown
	\item \AccountingIdle
	\item \AppStart~\textbf{TrickleApp}
	\item \ServerlessUp
	\item \AppRuns 
	\item \AccountingConsumes
	\item \UserSendsCredit~durante la ejecución y de forma periódica, poco ($ \approx 0.25~GRC$)
	\item \AppRuns~hasta su finalización
	\item \ServerlessDown
	\item \AccountingIdle
	\item \AccountingCredit~$\approx 0$
\end{itemize}

Sobre la aplicación \textbf{TrickleApp}:
\begin{itemize}
	\setlength\itemsep{0em}
	\item carga media (6-8 cores)
	\item tiene que ser estable
	\item puede ser algo sensible a cambios de recursos
	\item Candidata: alguna de procesamiento streaming
\end{itemize}

\newpage
\section{Credit-bounded processing}
\begin{itemize}
	\setlength\itemsep{0em}
	\item \ContainerUp
	\item \UserNoCredit
	\item \UserSendsCredit, suficiente ($\approx 10~GRC$)
	\item \AppStart~\textbf{GreedyApp}
	\item \ServerlessAdapts
	\item \AccountingConsumes
	\item \AppRuns~mientras tiene recursos, o sus métricas de rendimiento son buenas, hasta que se consume todo el crédito 
	\item \AccountingRestricts
	\item \AppRuns, pero ahora los recursos son mínimos, la aplicación decide parar
	\item \AccountingCredit~$\approx 0$, o incluso ligeramente negativo (deuda)
\end{itemize}

Sobre la aplicación \textbf{GreedyApp}:
\begin{itemize}
	\setlength\itemsep{0em}
	\item 
\end{itemize}

\newpage
\section{Credit-aware processing}
\begin{itemize}
	\setlength\itemsep{0em}
	\item \ContainerUp
	\item \UserNoCredit
	\item \UserAccountingActivated
	\item \AccountingRestricts
	\item \AppStart~\textbf{PoliteApp}
	\item \AppWaits~(ahora crédito $= 0~GRC$)
	\item \UserSendsCredit~($\approx 10~GRC$)
	\item \AccountingAllows
	\item \AppResume
	\item \ServerlessAdapts
	\item \AppRuns, hasta que el crédito esté próximo a agotarse, para no producir una restricción
	\item \AppWaits~(ahora crédito $\approx 0~GRC$)
	\item \ServerlessDown
	\item \AccountingIdle
	\item \Wait
	\item \UserSendsCredit~($\approx 10~GRC$)
	\item \AppResume, ahora que hay crédito disponible
	\item \AppRuns, hasta finalizar
	\item \ServerlessDown
	\item \AccountingIdle
	\item \AccountingCredit~$> 0$
\end{itemize}

Sobre la aplicación \textbf{PoliteApp}:
\begin{itemize}
	\setlength\itemsep{0em}
	\item 
\end{itemize}

\newpage
\section{Candidatos a aplicaciones}

\begin{itemize}
	\item Benchmarks de procesamiento de calculo
	\item Aplicaciones de metagenomica
\end{itemize}

\newpage

\section*{Acknowledgements}


\section*{References}

\bibliography{mybibfile}

\end{document}
